\begin{theorem}[DP32: sharp asymptotic slopes on the actual Honda dual]\label{DP32}
Let $p$ be any prime, $h\ge2$, $q=p^h$, $d=h-1$, and $\mathcal O=W(\F_q)$. The Cartier dual $H=G_h^D$ of the Honda group with logarithm $L_h(T)=\sum_{n\ge0}T^{q^n}/p^n$ has integral coordinates $X_1,\ldots,X_d$ and normalized primitive invariant derivations $D_i$ for which, setting
\[
 c_i=\frac{p^i}{q-1},\qquad r_i=\frac1{p-1}-c_i\quad(1\le i\le d),
\]
every integral formal function satisfies coefficientwise
\begin{equation}\label{eq:dual-mixed}
 v_p(D_1^{a_1}\cdots D_d^{a_d}f)
 \ge\max\left\{0,\ \sum_i c_i a_i-\frac{\sum_i s_p(a_i)}{p-1}\right\}.
\end{equation}
For every $i$ and $n=q^r$,
\begin{equation}\label{eq:dual-sharp}
 v_p(D_i^nX_i(0))=c_i(n-1).
\end{equation}
Thus every displayed directional slope is optimal, and $\sum_i c_i=c_h$. These assertions persist over $W(k)\supset\mathcal O$. In particular the order-$p$ obstruction DP30 does not imply absence of positive asymptotic gain on the dual.
\end{theorem}
\begin{proof}
We construct an integral group and prove its dual identification before estimating it. Over $\Z_p$, define a vector of monomials
\[
 M_1(\mathbf X)=X_d^{p^2},\qquad M_i(\mathbf X)=X_{i-1}^p\ (i\ge2),
 \qquad \mathcal L(\mathbf X)=\sum_{n\ge0}\frac{M^{\circ n}(\mathbf X)}{p^n}.
\]
For $d=1$, this means $M_1=X_1^{p^2}$. The degrees tend to infinity, so the series is well defined and has identity linear term. Its functional equations are
\[
 \mathcal L_1=X_1+p^{-1}\mathcal L_d(\mathbf X^{p^2}),\qquad
 \mathcal L_i=X_i+p^{-1}\mathcal L_{i-1}(\mathbf X^p)\ (i\ge2).
\]
Apply Hazewinkel \cite[Lemma 2.3, (2.2)--(2.6), p.~74]{Hazewinkel} with $B=\Z_p$, ideal $pB$, $\sigma=\id$, the lemma's $q=p$, matrix $s_1$ having $1/p$ just below the diagonal, and $s_2$ having $1/p$ in entry $(1,d)$. When $d=1$, $s_1=0$. It gives an integral commutative law $F_H=\mathcal L^{-1}(\mathcal L(\mathbf X)+\mathcal L(\mathbf Y))$ and an integral $p$-series. The exact identity
\[
 \mathcal L(M(\mathbf X))=p(\mathcal L(\mathbf X)-\mathbf X)
\]
and the congruence clause \cite[(2.8), p.~75]{Hazewinkel} imply $[p]_H\equiv M\pmod p$.

Here the finite-flat condition can be checked directly. Put $C=\Z_p[[U_1,\ldots,U_d]]$, $B=\Z_p[[X_1,\ldots,X_d]]$ and send $U_i$ to $[p]_{H,i}(\mathbf X)$. Modulo $p$, $B/p$ is free over $C/p$ on the $p^h$ monomials
\[
 X_1^{b_1}\cdots X_d^{b_d},\qquad
 0\le b_i<p\ (i<d),\quad 0\le b_d<p^2.
\]
Their lifts define $C^{p^h}\to B$. Successively lifting the coefficient expansion modulo $p$ and dividing the remainder by $p$ proves surjectivity by $p$-adic completeness. A relation has all coefficients divisible by $p$; dividing by $p$ is valid since $B$ is $p$-torsion-free. Repetition and $p$-adic separatedness prove injectivity. Thus $[p]_H$ is finite free of degree $p^h$. Its iterated kernels form a connected $p$-divisible group of height $h$ and dimension $d$, as in Tate \cite[\S2.2, Proposition 1, pp.~162--164]{Tate67}.

Everything is defined over $\Z_p$. On the special fibre over $\F_p$, a full cycle of $M$ has total Frobenius exponent $2+(d-1)=h$, so $[p]_H^d=M^d=F_H^h$. For $G=H^D$, dual dimensions give $\dim G=1$ \cite[\S2.3, Proposition 3, p.~166]{Tate67}, and duality gives $V_G^h=[p]_G^d$. Since $F_G^hV_G^h=[p]_G^h$ and $d=h-1$, cancellation gives $F_G^h=[p]_G$. Cancellation is valid because multiplication by $p^d$ is surjective as an fppf sheaf. All these maps are over $\F_p$, so they are actual endomorphisms with no unidentified twists. Their Frobenius kernels show that $G$ is connected. It therefore has an integral one-dimensional formal group over $\Z_p$ by Tate's equivalence. In an integral coordinate $T$ its $p$-series $f$ has linear term $pT$ and reduction $T^q$.

Over $\mathcal O=W(\F_q)$, this is precisely a Lubin--Tate series for uniformizer $p$. The known uniqueness and strict-isomorphism theorem \cite[\S5, pp.~76--77, (5.1)--(5.4)]{Hazewinkel} identifies the group with this $p$-series with the normalized Lubin--Tate group $F_f$, and identifies all such groups for the same uniformizer integrally and strictly. The Honda group $G_h$ has a $p$-series in this class. Thus $G\simeq G_h$ integrally over $\mathcal O$, and dualizing identifies the constructed $H$ with the actual Cartier dual claimed in the statement. No guessed dual logarithm or unproved lifting of an action is involved.

We now compute its radius in each coordinate. Put $S_h=(q-1)/(p-1)$, choose $\pi^{q-1}=p$, and set $\alpha_i=\pi^{S_h-p^i}$, so $v_p(\alpha_i)=r_i>0$. Scale coordinate $i$ in input and output by $\alpha_i$. For every arrow $X_j^{p^e}$ in $M_i$ one has $p^er_j-r_i=1$. Consequently the $n=1$ term of
\[
 \mathcal L'_i(\mathbf Z)=\alpha_i^{-1}\mathcal L_i(\alpha_1Z_1,\ldots,\alpha_dZ_d)
\]
has coefficient exactly one. If an $n$-fold cyclic iterate has successive exponents $e_1,\ldots,e_n$, its coefficient valuation after scaling is
\[
 1+p^{e_1}+p^{e_1+e_2}+\cdots+p^{e_1+\cdots+e_{n-1}}-n.
\]
This is positive for $n\ge2$. Hence $\mathcal L'$ is integral and reduces to $\mathbf Z+M(\mathbf Z)$. Its inverse $\mathcal E'$ is integral. In characteristic $p$, the monomial map $M$ is additive, and therefore
\[
 \overline{\mathcal E'}(\mathbf Z)=\sum_{n\ge0}(-1)^nM^{\circ n}(\mathbf Z).
\]
Composition with $\id+M$ verifies the cancellation, also at $p=2$. In component $i$, after setting other variables to zero, terms survive exactly at cyclic step numbers $dr$, with degrees $q^r$ and nonzero coefficients $(-1)^{dr}$. Undoing the input and output scales gives
\[
 v_p([t^{q^r}]\mathcal E_i(t e_i))=-(q^r-1)r_i.
\]

The normalized invariant derivations are integral and commute. In additive logarithm coordinates they are partial derivatives, so
\[
 D^{\mathbf a}f(0)=\mathbf a![\mathbf t^{\mathbf a}]f(\mathcal E(\mathbf t)).
\]
Integrality of $\mathcal E'$ says that every coefficient of $\mathcal E_j$ has weighted valuation at least $r_j$ when variable $t_i$ is given weight $r_i$. Integral substitution therefore bounds $[\mathbf t^{\mathbf a}]f(\mathcal E)$ by $-\sum_i r_i a_i$. Adding Legendre's exact factorial valuation proves \eqref{eq:dual-mixed} at zero; integrality also supplies the bound zero. Apply the same computation in $\mathbf Y$ to the integral translated function $f(F_H(\mathbf X,\mathbf Y))$, treating $\mathbf X$ as coefficient variables. Translation invariance gives the coefficientwise bound. At $n=q^r$, the exact exponential coefficient above and $s_p(n)=1$ prove \eqref{eq:dual-sharp}. Finally
\[
 \sum_{i=1}^{h-1}c_i=\frac{p+\cdots+p^{h-1}}{p^h-1}=c_h.
\]
This proves all assertions, including sharpness against a larger slope with sublinear error.
\end{proof}
